\section{Limitations}
\label{sec:limitations}

A substrate-layer construction has structural limits that the section above implicitly assumed; this section names them.

\paragraph{Adoption is the binding constraint.}
A substrate-layer guarantee holds only for deployments that route every tool call through the admission hook. Today, most agentic-AI deployments do not. The Model Context Protocol~\cite{anthropicMCP2024} and its analogues are passive dispatchers in their current implementations. Adopting admission-layer discipline requires the deployer to install an admission hook, classify each tool's methods, implement rollback executors for reversible-class methods, and route every call through the hook without side channels. None of these tasks is research; all of them are work. The paper's structural argument is independent of how widely the discipline is adopted, but the deployed safety property obviously is not. We are honest about this: a substrate-layer guarantee is a structural construction, and its operational value is bounded by the deployer's discipline in routing calls through the substrate.

\paragraph{Class registries are operator-trusted.}
The grammar pushes the classification of methods into the operator's class registry. A method whose effects are destructive but is mis-classified as reversible admits without bilateral cosignature; the substrate's discipline does not catch this. The operator's classification is, in a precise sense, the substrate's trust assumption. Strengthening this position with formal-methods discipline on the registry (a typed schema for the registry, a class-derivation tool, a per-method test that empirically verifies reversibility) is downstream work; this paper does not discharge it.

\paragraph{Reversibility is empirical, not formal.}
A rollback executor for a reversible-class method discharges the rollback obligation only if it actually inverts the method's effect. In practice, this is an empirical claim about the executor's correctness against the target system. A file-quarantine rollback that fails because the file's parent directory was removed in the interim discharges nothing; the rollback executor returns an error. The substrate's response to such failures is the failed-rollback amendment obligation~\cite{programmableSovereignty2026}, which converts the rollback failure into a publicly-witnessed event rather than a silent retry, but the substrate cannot synthesise a working rollback from a broken one. Reversibility is a property of the method-plus-environment, not of the typed envelope.

\paragraph{Destructive-class calls happen and are sometimes correct.}
The grammar does not refuse destructive calls. It admits them only under bilateral cosignature. A correctly-classified destructive call that both polities admit will execute; the substrate's discipline protects the case where one polity admits and the other does not, not the case where both agree. A substrate-layer construction cannot guarantee that the right destructive calls happen and the wrong ones do not; it guarantees only that destructive calls admitted by a single party do not execute. This is a weaker claim than ``the substrate enforces correctness'', and it is the strongest claim the construction supports.

\paragraph{The threat model is bounded.}
The threat model in Section~\ref{sec:threat} covers a misaligned model, strategic composition, alignment-faking, and operator manipulation. It does not cover an adversary who has compromised the substrate's runtime, who can rewrite the class registry, who can forge cosignatures, or who can bypass the admission hook entirely. Substrate-level adversaries fall outside the grammar; defending against them is a substrate-correctness problem, not a substrate-grammar problem.

\paragraph{No empirical evaluation.}
This paper does not measure. The companion substrate paper~\cite{programmableSovereignty2026} carries a corpus evaluation in the endpoint-detection-and-response deployment instance; the present paper is a position paper that argues for the construction's value in the agentic-AI setting. A targeted empirical evaluation, for example a benchmark of misaligned agent behaviours against substrate-admission-protected and unprotected deployments, is the natural follow-up. We expect that an evaluation will sharpen the construction (by exposing classification mistakes, executor bugs, and protocol-bypass paths) more than it will validate the headline claim. The structural argument stands or falls on its own merits; the empirical work is the next step.

\paragraph{Scope of the formal claim.}
The headline structural theorems live in the companion substrate~\cite{programmableSovereignty2026}, not in this paper. The present paper proposes a reframing of agentic-AI safety as a layered construction; the load-bearing formal content is cited rather than developed here. Reviewers who want machine-checked claims about admission, capability attenuation, and amendment refinement should read the companion paper. Reviewers who want the structural argument for why such claims matter in the agentic-AI setting will find that argument here.
